% Preámbulo común del material docente · XeLaTeX
\documentclass[11pt,a4paper]{article}
\usepackage[margin=22mm,top=20mm,bottom=22mm]{geometry}
\usepackage{fontspec}
\setmainfont{Avenir Next}[BoldFont={Avenir Next Demi Bold}, BoldItalicFont={Avenir Next Demi Bold Italic}]
\setsansfont{Avenir Next}
\setmonofont{Menlo}[Scale=0.88]
\newfontfamily\symfont{Apple Symbols}
\newfontfamily\unifont{Arial Unicode MS}
\usepackage{unicode-math}
\setmathfont{latinmodern-math.otf}
\usepackage{polyglossia}\setdefaultlanguage{spanish}
\usepackage{xcolor,booktabs,tabularx,enumitem,parskip,microtype,titlesec,fancyhdr}
\usepackage[most]{tcolorbox}
\usepackage[hidelinks]{hyperref}
\emergencystretch=2em

\definecolor{tinta}{HTML}{1F2933}
\definecolor{acento}{HTML}{B7410E}
\definecolor{suave}{HTML}{F4F1EC}
\definecolor{gris}{HTML}{6B7280}
\definecolor{linea}{HTML}{D6D0C6}
\color{tinta}

\titleformat{\section}{\Large\bfseries\color{acento}}{\thesection}{0.6em}{}
\titleformat{\subsection}{\large\bfseries}{\thesubsection}{0.6em}{}
\titlespacing*{\section}{0pt}{1.6em}{0.5em}
\titlespacing*{\subsection}{0pt}{1.1em}{0.3em}
\setlist{nosep, leftmargin=1.4em, itemsep=2pt}
\renewcommand{\arraystretch}{1.25}

% bloque de órdenes de terminal
\newtcblisting{orden}{listing only, colback=suave, colframe=linea, boxrule=0.4pt, arc=2pt,
  left=6pt, right=6pt, top=4pt, bottom=4pt, listing options={basicstyle=\ttfamily\small, breaklines=true, columns=fullflexible, keepspaces=true}}
% nota destacada
\newtcolorbox{nota}[1][Nota]{colback=white, colframe=acento, boxrule=0.6pt, arc=2pt, title=#1,
  fonttitle=\bfseries\small, coltitle=white, colbacktitle=acento, left=6pt, right=6pt}
% preguntas
\newcommand{\pregunta}[2]{\item[\textbf{(#1)}] #2}
\newenvironment{preguntas}{\begin{description}[leftmargin=2.4em, style=nextline, itemsep=4pt]}{\end{description}}
\newcommand{\ok}{{\symfont ✔}}
\newcommand{\nok}{{\symfont ✘}}
\newcommand{\qed}{{\symfont ∎}}
\newcommand{\codigo}[1]{\texttt{#1}}

\pagestyle{fancy}\fancyhf{}
\renewcommand{\headrulewidth}{0pt}
\fancyfoot[C]{\small\color{gris}\docpie\ · página \thepage}

\newcommand{\cabecera}[3]{%
  {\Huge\bfseries\color{acento}#1\par}\vspace{4pt}
  {\large #2\par}\vspace{2pt}
  {\color{gris}#3\par}\vspace{8pt}
  {\color{linea}\hrule height 0.6pt}\vspace{10pt}}

\newcommand{\docpie}{Laboratorio 2 · guía del estudiante}
\begin{document}
\cabecera{Laboratorio 2 · Controlar una iteración}{Un agente que ajusta el paso mirando solo dos números, y el teorema de Banach}{Métodos numéricos · Matemáticas · 90 minutos · en parejas}

En el laboratorio 1 el agente aprendió \emph{dónde cortar}. Hoy aprende algo más difícil: \textbf{ajustar un parámetro sobre la marcha} mirando solo lo que la iteración le deja ver. Y el teorema es el de Banach, el que está detrás de casi todos los métodos iterativos del curso.

Todas las órdenes se escriben dentro de \codigo{02\_punto\_fijo} y están escritas para Windows.
\begin{nota}[Si usas Mac o Linux]
\begin{center}
\begin{tabular}{ll}
\textbf{Windows} & \textbf{Mac / Linux} \\
\codigo{.venv\textbackslash Scripts\textbackslash python} & \codigo{.venv/bin/python} \\
\codigo{..\textbackslash docs\textbackslash curso\textbackslash lab02\textbackslash experimentos.py} & \codigo{../docs/curso/lab02/experimentos.py} \\
\end{tabular}
\end{center}
\end{nota}

\subsection*{Qué vas a hacer}
\begin{enumerate}
\item Ver que un parámetro fijo no sirve, y que un agente que solo mira dos números lo ajusta bien.
\item Quitarle al agente información y ver qué deja de poder aprender (\textbf{el estado importa}).
\item Darle una recompensa de progreso que él mismo puede manipular (\textbf{el potencial importa}).
\item Dibujar el grafo del teorema de Banach y encontrar el hueco cuando falta la contracción.
\end{enumerate}

\section{Instalación \hfill\small\normalfont 0–10 min}
\begin{orden}
cd 02_punto_fijo
python -m venv .venv
.venv\Scripts\python -m pip install -e ".[dev]"
.venv\Scripts\python -m fixpointrl --no-tui --seed 1
\end{orden}
Cada proyecto tiene su propio \codigo{.venv}: el de \codigo{01\_biseccion} no sirve aquí.

\section{El problema \hfill\small\normalfont 10–20 min, en papel}

Para resolver $f(x)=0$ iteramos $x \leftarrow g(x)$ con $g(x) = x - \alpha f(x)$. La raíz $r$ es punto fijo de $g$ \textbf{para cualquier $\alpha$}. Pero la iteración solo converge si $g$ es contracción cerca de $r$:
\[
|g'(r)| = |1 - \alpha f'(r)| < 1 .
\]
El agente \textbf{no conoce $f'$ ni $r$}. Tras cada iteración solo ve dos cosas medibles:
\begin{itemize}
\item $\rho = |f(x_{n+1})| / |f(x_n)|$, la razón de residuos consecutivos (en 6 cubos);
\item si el residuo \textbf{cambió de signo} (oscila) o no (monótona).
\end{itemize}
Y puede: mantener $\alpha$, doblarlo, reducirlo a la mitad o invertir su signo.

\begin{nota}[Responde antes de correr nada]
\begin{preguntas}
\pregunta{a}{Para $f(x) = 2(x-1)$, ¿para qué valores de $\alpha$ converge la iteración desde cualquier $x_0$? ¿Y qué $\alpha$ la hace converger en un solo paso?}
\pregunta{b}{Cerca de la raíz, $\rho \approx |1-\alpha f'(r)|$. Di qué harías tú: si $\rho>1$ y la iteración es monótona, ¿qué está mal? ¿Y si $\rho>1$ y oscila? ¿Y si $\rho<1$ pero oscila?}
\pregunta{c}{Predice: ¿en cuál de los 12 estados (6 cubos de $\rho$ $\times$ monótona/oscila) es más difícil decidir?}
\end{preguntas}
\end{nota}

\section[Sin agente: alfa fijo]{Sin agente: $\alpha$ fijo \hfill\small\normalfont 20–30 min}
\begin{orden}
.venv\Scripts\python ..\docs\curso\lab02\experimentos.py alfa_fijo
\end{orden}
\begin{preguntas}
\pregunta{d}{¿Cuál es el mejor $\alpha$ fijo y en qué porcentaje de problemas converge? ¿Por qué ningún $\alpha$ fijo puede funcionar siempre? (Pista: el signo de $f'(r)$ cambia de un problema a otro.)}
\end{preguntas}

\section{Mira cómo aprende \hfill\small\normalfont 30–45 min}
\begin{orden}
.venv\Scripts\python -m fixpointrl --seed 1
\end{orden}
Anota la generación de cada hito. Luego corre con \codigo{--seed 2} y completa la segunda columna.

\renewcommand{\arraystretch}{2.1}
\begin{tabularx}{\linewidth}{|X|c|c|}
\hline
\textbf{hito} & \textbf{gen. (seed 1)} & \textbf{gen. (seed 2)} \\
\hline
«si diverge monótonamente, el signo de $\alpha$ está mal» & \hspace{2.6cm} & \hspace{2.6cm} \\
\hline
primera convergencia & & \\
\hline
«si oscila, $\alpha$ es demasiado grande» & & \\
\hline
ley coincide con Banach en los 12 estados & & \\
\hline
primera demostración mínima & & \\
\hline
\end{tabularx}
\renewcommand{\arraystretch}{1.25}

Abre \codigo{runs\textbackslash <fecha>\textbackslash informe.md} y mira la tabla \textbf{«ley de control aprendida»}.
\begin{preguntas}
\pregunta{e}{Compárala con tu tabla de (b). ¿En qué estados coincide con lo que tú habrías hecho?}
\pregunta{f}{En el estado «$\rho \ge 1{,}5$, monótona» la teoría admite dos acciones (invertir o reducir). ¿Por qué las dos son razonables ahí?}
\end{preguntas}

\section{El estado importa \hfill\small\normalfont 45–60 min}
Le quitamos al agente el bit de oscilación: ahora solo ve $\rho$.
\begin{orden}
.venv\Scripts\python ..\docs\curso\lab02\experimentos.py ciego
\end{orden}
\begin{preguntas}
\pregunta{g}{¿Cuánto cae el éxito? En los estados con $\rho>1$, ¿qué acción aprende y por qué es la mejor que puede tomar \textbf{sin saber} si oscila?}
\pregunta{h}{Propón otro observable (algo que se pueda medir sin conocer $f'$ ni $r$) que le ayudaría más que el bit de oscilación, o argumenta que no hace falta.}
\end{preguntas}

\section{El potencial importa \hfill\small\normalfont 60–72 min}
La recompensa de progreso es $\log_{10}\bigl(|f(x_n)|/|f(x_{n+1})|\bigr)$: premia reducir el residuo, que el agente no controla directamente. Vamos a cambiarla por $\log_{10}\bigl(|\text{paso anterior}|/|\text{paso nuevo}|\bigr)$: premia \textbf{moverse cada vez menos}.
\begin{orden}
.venv\Scripts\python ..\docs\curso\lab02\experimentos.py potencial
\end{orden}
\begin{preguntas}
\pregunta{i}{¿Cuánto cae el éxito? Mira los estados marcados \nok: ¿qué hace ahora el agente cuando oscila cerca de la raíz ($\rho<0{,}25$)? ¿Por qué «moverse menos» le da puntos aunque no le acerque a la raíz? ¿Con qué acción se consigue moverse menos sin resolver nada?}
\end{preguntas}

\section{El teorema de Banach como grafo \hfill\small\normalfont 72–85 min}
Abre \codigo{fixpointrl\textbackslash proof\_kb.py}: 13 pasos con dependencias y 7 distractores.
\begin{preguntas}
\pregunta{j}{Dibuja el grafo en papel (nodos, flechas, hipótesis, \qed). Luego compáralo con el del programa:}
\end{preguntas}
\begin{orden}
.venv\Scripts\python ..\docs\curso\lab02\experimentos.py dibujar
\end{orden}
\begin{preguntas}
\pregunta{k}{El distractor \codigo{D\_LOCAL} dice «$|g'(p)|<1$, luego converge desde cualquier $x_0$». ¿Qué diferencia hay entre lo que dice el teorema y lo que dice el distractor? (Pista: local frente a global; ¿qué hipótesis garantiza que la iteración no se sale de $[a,b]$?)}
\end{preguntas}
Ahora el hueco:
\begin{orden}
.venv\Scripts\python ..\docs\curso\lab02\experimentos.py grafo
\end{orden}
Quitamos H2 (la contracción) de las dependencias de \codigo{ERR}. El verificador acepta una demostración donde la contracción \textbf{se declara después de haberla usado}.
\begin{preguntas}
\pregunta{l}{\codigo{ERR} dice $|x_{n+1}-p| \le k\,|x_n-p|$. ¿De dónde sale esa $k$? ¿Qué le pasa a la convergencia (\codigo{LIM}) si solo sabemos $k\le 1$? Relaciónalo con el distractor \codigo{D\_K1}.}
\pregunta{m}{¿Quién es responsable de que la demostración sea correcta: el agente, el verificador o quien escribió las \codigo{deps}?}
\end{preguntas}

\section{Cierre \hfill\small\normalfont 85–90 min}
Entrega (una por pareja): la tabla de hitos, las respuestas (a)–(m) en la hoja de entrega, y el párrafo final: \emph{¿qué observable elegirías para controlar el paso de Newton (laboratorio 4) sin conocer $f''$?}

\begin{nota}[Tarea opcional]
El estado tiene 12 combinaciones. Cambia \codigo{RATIO\_EDGES} en \codigo{env\_fixpoint.py} a solo dos cortes, \codigo{(0.5, 1.0)}, corre \codigo{--no-tui} y compara el éxito. ¿Cuánta resolución en $\rho$ hace falta? Revierte el cambio después.
\end{nota}

\clearpage
\cabecera{Hoja de entrega · Laboratorio 2}{Nombres: \rule{7cm}{0.4pt}}{Fecha: \rule{3cm}{0.4pt} \quad Semillas usadas: \rule{2cm}{0.4pt}}

\begin{respuesta}[1.9cm]{(a) $\alpha$ que convergen para $f(x)=2(x-1)$, y el que converge en un paso}\end{respuesta}
\begin{respuesta}[2.6cm]{(b) qué harías en cada caso: $\rho>1$ monótona · $\rho>1$ oscila · $\rho<1$ oscila}\end{respuesta}
\begin{respuesta}[1.6cm]{(c) el estado más difícil de decidir}\end{respuesta}
\begin{respuesta}[2.2cm]{(d) mejor $\alpha$ fijo, su porcentaje, y por qué ninguno sirve siempre}\end{respuesta}
\begin{respuesta}[1.9cm]{(e) estados donde la ley aprendida coincide con tu tabla}\end{respuesta}
\begin{respuesta}[1.9cm]{(f) por qué invertir y reducir son ambas razonables con $\rho\ge 1{,}5$ monótona}\end{respuesta}
\begin{respuesta}[2.8cm]{(g) estado ciego: cuánto cae el éxito, qué aprende con $\rho>1$ y por qué}\end{respuesta}
\begin{respuesta}[2.4cm]{(h) otro observable, o por qué no hace falta}\end{respuesta}

\clearpage
\begin{respuesta}[3.4cm]{(i) potencial tramposo: cuánto cae, qué hace en $\rho<0{,}25$ oscila, con qué acción se «mueve menos»}\end{respuesta}
\begin{respuesta}[7.5cm]{(j) el grafo del teorema de Banach}\end{respuesta}
\begin{respuesta}[2.4cm]{(k) teorema (global) frente a \codigo{D\_LOCAL}: qué hipótesis falta}\end{respuesta}
\begin{respuesta}[2.8cm]{(l) de dónde sale $k$; qué pasa con \codigo{LIM} si $k\le 1$; relación con \codigo{D\_K1}}\end{respuesta}
\begin{respuesta}[1.6cm]{(m) quién es responsable}\end{respuesta}
\begin{respuesta}[3.0cm]{Párrafo final: un observable para controlar el paso de Newton sin conocer $f''$}\end{respuesta}
\end{document}
